\subsection*{The DMA Statement Family}

All of Croqtile's tile-movement operations use the same basic form:

\begin{lstlisting}[language=croqtile,caption={The general DMA form}]
future = engine.operation.modifiers source => destination;
\end{lstlisting}

The important pieces are:

\begin{table}[htbp]
\centering\small
\begin{tabularx}{\linewidth}{@{}lX@{}}
\toprule
\textbf{Part} & \textbf{Meaning} \\
\midrule
\texttt{future} & Optional handle to the transfer result. Use \texttt{future.data} to access the moved tile and \texttt{future.span} to query its shape. \\
\texttt{operation} & \texttt{copy}, \texttt{transp\textless{}.\allowbreak{}.\allowbreak{}.\allowbreak{}\textgreater{}}, or \texttt{pad\textless{}.\allowbreak{}.\allowbreak{}.\allowbreak{}\textgreater{}}. \\
\texttt{modifiers} & Optional flags such as \texttt{.async}, \texttt{.zfill}, or \texttt{.\allowbreak{}swiz\textless{}N\textgreater{}}. \\
\texttt{source =\textgreater{} destination} & The source view and the destination memory or destination view. \\
\bottomrule
\end{tabularx}
\end{table}

Minimal examples:

\begin{lstlisting}[language=croqtile,caption={A concrete copy}]
f = dma.copy input.chunkat(block) => shared;
g = dma.transp<1, 0> matrix => local;
h = dma.pad<{2, 1}, {3, 2}, {0, 0}, 0> tile => shared;
\end{lstlisting}

The key shift from Section~\ref{sec:1-7} is that these statements move shaped regions, not individual elements.
